\section{Methodology}
\label{sec:methodology}

\subsection{System Architecture}
The report uses a version-focused methodology. Rather than treating the repository as a single static compressor, we identify three representative designs. \texttt{v5} is the main balanced design, combining classic modeling/coding separation with practical runtime. \texttt{v9} is the stronger compression-oriented extension of the same family. \texttt{v8} is the speed-focused contrast that removes the statistical pipeline and uses direct dictionary matching. This framing exposes the central trade-off between ratio, throughput, and overall balance.

\subsection{Algorithm Selection}
The \texttt{v5} line remains the primary subject because it offers the best overall speed/ratio trade-off in the refreshed benchmark set and best illustrates the repository's statistical compression pipeline. BWT, MTF, and zero-run coding reshape low-entropy structure before modeling; order-1 adaptive modeling captures local byte dependencies better than a static order-0 model; and range coding/rANS provide the coding stage while keeping the implementation efficient. \texttt{v9} is included because it improves ratio further through per-block candidate selection and broader transform choices, while \texttt{v8} is kept as a comparison point because its strength is speed, not compactness.

\subsection{Optimization Strategy}
Development proceeded iteratively. Earlier versions established the basic order-0 and range-coding path, then introduced BWT preprocessing and faster data structures. The \texttt{v5} line consolidated the strongest ideas into one pipeline: block-based BWT, MTF, optional zero-run encoding, adaptive order-1 modeling, and parallel block processing. \texttt{v8} then explored the opposite extreme by discarding preprocessing and entropy coding entirely in order to maximize speed. Finally, \texttt{v9} extended the statistical line with per-block candidate search over multiple transforms and model orders to push ratio further.

\subsection{Benchmarking}
All claims in this report are based on the benchmark corpus used throughout the project. The corpus contains eight files (A--H) totaling 13 MiB, with different characteristics including text, structured binaries, high-entropy data, and an incompressible random case. The reported metrics are compressed size, compression ratio, bits per symbol, compression time, decompression time, and lossless correctness. The benchmarks compare \texttt{v5}, \texttt{v8}, and \texttt{v9} against gzip, bzip2, xz, and zstd.
Benchmarks used the standard release configuration from the project Makefile (\texttt{g++}/C++17 with \texttt{-O3 -march=native -flto=auto}) on the same 8-file, 13 MiB corpus.
